% ---
% id: protocol-compiler/thm:sifting-and-key
% kind: thm
% title: "Announcement locality, sifting and party computability"
% summary: >-
%   Each announcement is a function of its sender's readout and of the
%   announcements before it, sifting is determined by the announcements, and
%   the holder computes the raw key locally. An ensemble-form measurement slot
%   of Alice's denotes a POVM.
% prominence: standard
% status: accepted
% uses: [def:compile, def:declaration, def:target, def:spec, def:well-formed-spec, def:exact]
% assumes: [A-SPEC]
% refs: []
% provenance: newly-derived
% implements: certiqs_key.protocol_compiler.spec:derived_holder_key_rule
% lean:
% notes: "Part (v) is proved here rather than imported, so that this contract is self-contained; meat-core/lem:ensemble-form proves the same. Only A-SPEC (ii) is spent, at (ii) and (vi); A-ALICE-MEAS and A-ENS are not declared, no step reading the measurement Alice performs or the physical source ensemble."
% ---

\begin{theorem}\label{thm:sifting-and-key}
Let $\Spec$ be a \tref{well-formed specification}
(Definition~\ref{def:well-formed-spec}) satisfying
Assumption~\ref{ass:spec}, so
that in particular $r^{0}_{\hold}$ maps into $\Zval$
(Assumption~\ref{ass:spec}\,(ii)), and let $\Decl$ be emitted by
Definition~\ref{def:compile}. Then:
\begin{enumerate}[label=(\roman*),leftmargin=2.2em]
\item \emph{(announcement locality)} $a_A$ is a function of
$\rd{\mathsf A}(x,v)$ alone and $a_B$ of $\rd{\mathsf B}(x,v)$ alone, and the
joint announcement is \eqref{eq:jointann};
\item \emph{(D(v))} the specification's $\varsigma$ --- consumed, and not
itself a slot of \eqref{eq:declaration} --- witnesses announcement-determined
sifting: $r(x,v)=\perp\iff\varsigma(a(x,v))=\mathrm{fail}$;
\item \emph{(D(viii))} $r$ is party-computable: the total map
\begin{equation}\label{eq:rhold}
  r_{\hold}(o,\kappa)\coloneqq
  \begin{cases}
    r^{0}_{\hold}(o,\kappa), & \varsigma(\kappa)=\mathrm{pass},\\
    \perp, & \varsigma(\kappa)=\mathrm{fail},
  \end{cases}
  \qquad (o,\kappa)\in\Oset{\hold}\times\Kann,
\end{equation}
satisfies \eqref{eq:partycomp}.
\item \emph{(\tref[transcript protocol]{transcript protocols})} if the
announcement slot holds a transcript protocol
(Definition~\ref{def:declaration}\,(DX5)), then for
every $m$ the $m$-th component of $a(x,v)$ is a function of
$\rd{\pi_m}(x,v)$ and of the components before it, and of nothing else; $a$ is
total; and (ii) and (iii) hold with this $a$. For $L=2$,
$\pi=(\mathsf A,\mathsf B)$ and $a_2$ independent of $\kappa_1$ this is (i);
\item \emph{(\tref{ensemble form})} if Alice's measurement slot is declared
in ensemble form and the conditions of Definition~\ref{def:exact} for that
form hold, then the operators \eqref{eq:ensemblepovm} satisfy
$M^{(t)}_x\succeq0$ and $\sum_xM^{(t)}_x=\idon{A_t}$ for every $t$, and
$(\tau^{(t)}_A)^{1/2}M^{(t)}_x(\tau^{(t)}_A)^{1/2}=N^{(t)}_x$: the denoted
declaration satisfies the POVM clause of D(iii) on $A$ as an identity of real
numbers, and $\tr[M_x\tau_A]=\sum_t\tr N^{(t)}_x$.
\item \emph{(the compiled general format)} if $\Spec$ is presented as
Definition~\ref{def:spec}\,(a$'$)--(e$'$) and $\Decl$ is emitted by the
general-format clauses of Definition~\ref{def:compile}, then: the
announcement slot of $\Decl$ holds the transcript $\Spec_{\mathrm{pub}}$
declares, message for message; the joint announcement function of $\Decl$ is
\eqref{eq:transcript} evaluated at \emph{those} tables and is a slot of
nothing; and Alice's measurement slot holds the form the transmitter artefact
declares. Consequently (iv) applies to $\Decl$, hence (ii) and (iii) hold of
it, and (v) applies to it exactly where that form is the ensemble form.
\end{enumerate}
Moreover, under the \tref{default party partition} these read:
$a_A:\Xset\to\CAset$, $a_B:\Vset\to\CBset$, and $r(x,v)=r_{\mathsf A}(x,a(x,v))$
for $\hold=\mathsf A$, respectively $r(x,v)=r_{\mathsf B}(v,a(x,v))$ for
$\hold=\mathsf B$.
\end{theorem}

\begin{proof}
(i) is Definition~\ref{def:spec}\,(d) together with \eqref{eq:jointann}: the
\tref[announcement map]{announcement maps} are declared on the readout alphabets, so no other
signature is expressible.

(ii) Fix $(x,v)$ and write $\kappa=a(x,v)$. If
$\varsigma(\kappa)=\mathrm{fail}$ then \eqref{eq:rdef} gives $r(x,v)=\perp$. If
$\varsigma(\kappa)=\mathrm{pass}$ then
$r(x,v)=r^{0}_{\hold}(\rd{\hold}(x,v),\kappa)\in\Zval$, and $\perp\notin\Zval$
by $\Zset=\Zval\sqcup\{\perp\}$, so $r(x,v)\ne\perp$. The two implications are
the two directions of the equivalence.

(iii) $r_{\hold}$ of \eqref{eq:rhold} is total on $\Oset{\hold}\times\Kann$
because $\varsigma$ is total and both branches are defined everywhere on their
respective preimages. Fix $(x,v)$ and put $o=\rd{\hold}(x,v)$,
$\kappa=a(x,v)$. If $\varsigma(\kappa)=\mathrm{pass}$, both sides of
\eqref{eq:partycomp} equal $r^{0}_{\hold}(o,\kappa)$ by \eqref{eq:rdef} and
\eqref{eq:rhold}; if $\varsigma(\kappa)=\mathrm{fail}$, both equal $\perp$.

(iv) By \eqref{eq:transcript} the $m$-th component is
$a_m(\rd{\pi_m}(x,v),\kappa_1(x,v),\dots,\kappa_{m-1}(x,v))$, and $a_m$ has
no other argument, its domain being declared as
$\Oset{\pi_m}\times\mathcal{C}_1\times\dots\times\mathcal{C}_{m-1}$; totality
follows by induction on $m$ from totality of the tables and of the
\tref[readout map]{readout maps}. The proofs of (ii) and (iii) use of $a$ only that it is a
total map into $\Kann$ on which $\varsigma$ and $r^{0}_{\hold}$ are defined, so
they apply unchanged. With $L=2$ and $a_2(o,\kappa_1)=a_B(o)$,
\eqref{eq:transcript} is \eqref{eq:jointann}.

(v) Fix $t$ and drop it. $\tau_A\ne0$ is positive semidefinite, so
$\tr\tau_A>0$ and $w_x\coloneqq\tr N_x/\tr\tau_A\ge0$ with $\sum_xw_x=1$ by
$\sum_xN_x=\tau_A$. With $S\coloneqq\tau_A^{-1/2}$ on the support, zero off
it, and $P$ the support projector: $SN_xS\succeq0$ and
$w_x(\idon{}-P)\succeq0$, so $M_x\succeq0$, and
$\sum_xM_x=S\tau_AS+(\idon{}-P)=\idon{}$. From
$0\preceq N_x\preceq\tau_A$, every vector in the kernel of $\tau_A$ is in
the kernel of $N_x$, so $PN_xP=N_x$; then
$\tau_A^{1/2}M_x\tau_A^{1/2}=PN_xP+w_x\tau_A^{1/2}(\idon{}-P)\tau_A^{1/2}
=N_x$, using $\tau_A^{1/2}S=P$ and $\tau_A^{1/2}(\idon{}-P)=0$. The trace
statement follows by taking traces and summing over $t$.

(vi) The three clauses are what the general-format clauses of
Definition~\ref{def:compile} do, and they are read off that definition and
not inferred: the transcript is passed through verbatim, \emph{$a$ is not a
slot} but the composition \eqref{eq:transcript} of the tables passed through,
and Alice's measurement slot is passed through in the form the artefact
declares, that definition converting neither form into the other. For the
consequences: Assumption~\ref{ass:spec}\,(ii) gives $r^{0}_{\hold}$ into
$\Zval$ and Definition~\ref{def:well-formed-spec}\,(iii) the totality of the
declared public tables, which is all (iv) consumes
besides the signatures, and the signatures are the ones (d$'$) of
Definition~\ref{def:spec} declares; so (iv) applies, and (ii), (iii) with it,
by (iv)'s last two sentences. Where the slot is the ensemble slot, the
conditions (v) is stated under are exactly the general-format conditions
Definition~\ref{def:exact} states for that form --- (S), (H), (P) of
$\{N^{(t)}_x\}_x$ per Alice sector, the identity of its declared-order sum
with $\tau^{(t)}_A$, and $\tau^{(t)}_A\ne0$ --- so (v) applies to it.
\emph{Nothing of (i) is claimed of such a declaration}: (i) is the two-map
case, which is $L=2$ with $a_2$ independent of $\kappa_1$, and (iv)'s last
sentence is where the two meet.

The default-partition statements follow by substituting
$\rd{\mathsf A}(x,v)=x$ and $\rd{\mathsf B}(x,v)=v$.
\end{proof}
